% Part: first-order-logic
% Chapter: proof-systems
% Section: natural-deduction

\documentclass[../../../include/open-logic-section]{subfiles}

\begin{document}

\iftag{FOL}
      {\olfileid{fol}{prf}{ntd}}
      {\olfileid{pl}{prf}{ntd}}

\olsection{استنتاج طبیعی}

استنتاج طبیعی دستگاهی برای !!{derivation} است که بناست استدلال واقعی
(به‌ویژه آن نوع استدلال منضبطی که ریاضی‌دانان به کار می‌برند) را بازتاب دهد.
استدلال واقعی از شماری الگوی «طبیعی» پیروی می‌کند. برای نمونه، برهان بر حسب
حالات به ما امکان می‌دهد نتیجه‌ای را بر پایهٔ مقدمه‌ای فصلی اثبات کنیم، به این
شیوه که نشان دهیم نتیجه از هر یک از دو طرف فصل به دست می‌آید. برهان غیرمستقیم
به ما امکان می‌دهد نتیجه‌ای را با نشان دادن اینکه نقیض آن به تناقض می‌انجامد
اثبات کنیم. برهان شرطی گزاره‌ای شرطی به صورت «اگر \dots آنگاه \dots» را با
نشان دادن اینکه تالی از مقدم به دست می‌آید اثبات می‌کند. استنتاج طبیعی
صورت‌بندی صوری برخی از این استنتاج‌های طبیعی است. هر یک از رابط‌های منطقی و
سورها دو قاعده دارد، یک قاعدهٔ معرفی و یک قاعدهٔ حذف، و هر کدام با یکی از
این الگوهای طبیعی استنتاج متناظر است. برای نمونه، $\Intro{\lif}$ با برهان شرطی
و $\Elim{\lor}$ با برهان بر حسب حالات متناظر است. قاعده‌ای به‌ویژه ساده
$\Elim{\land}$ است که از $!A \land !B$ استنتاج~$!A$ (یا $!B$) را مجاز می‌سازد.

یکی از ویژگی‌هایی که استنتاج طبیعی را از دیگر دستگاه‌های !!{derivation} متمایز
می‌کند کاربرد فرض‌هاست. در استنتاج طبیعی، هر !!^a{derivation} درختی از
!!{formula}s است. یک !!{formula} در ریشهٔ درختِ !!{formula}s قرار دارد و
«برگ‌های» درخت همان !!{formula}sیی هستند که نتیجه از آن‌ها اشتقاق می‌یابد.
در استنتاج طبیعی، بعضی !!{formula}sی برگ درون !!{derivation} نقشی دارند، اما
تا هنگامی که !!{derivation} به نتیجه برسد «مصرف شده‌اند». این با آن رویه در
استدلال واقعی متناظر است که در آن فرضیه‌هایی را وارد می‌کنیم که فقط مدتی کوتاه
معتبر می‌مانند. برای نمونه، در برهان بر حسب حالات، صدق هر یک از دو طرف فصل را
فرض می‌کنیم؛ در برهان شرطی، صدق مقدم را فرض می‌کنیم؛ و در برهان غیرمستقیم،
صدق نقیض نتیجه را فرض می‌کنیم. وارد کردن فرض‌های موقت به این شیوه و سپس کنار
گذاشتن آن‌ها برای اثبات یک گام میانی، مشخصهٔ استنتاج طبیعی است. فرمول‌های
موجود در برگ‌های یک !!{derivation} در استنتاج طبیعی «فرض» نامیده می‌شوند و برخی
از قواعد استنتاج می‌توانند آن‌ها را «!!{discharge}» کنند. برای نمونه، اگر
!!^a{derivation}ی برای~$!B$ از فرض‌هایی داشته باشیم که~$!A$ را نیز دربر می‌گیرند،
قاعدهٔ $\Intro{\lif}$ به ما امکان می‌دهد~$!A \lif !B$ را استنتاج و هر فرضی به
صورت~$!A$ را رفع کنیم. (برای ثبت اینکه کدام فرض‌ها در کدام استنتاج‌ها رفع
می‌شوند، استنتاج و فرض‌هایی را که رفع می‌کند با یک عدد نشانه‌گذاری می‌کنیم.)
فرض‌های !!{undischarged} که در پایان !!{derivation} می‌مانند، در مجموع
برای صدق نتیجه کافی‌اند؛ بنابراین هر !!^a{derivation} ثابت می‌کند که فرض‌های
!!{undischarged} آن مستلزم نتیجهٔ آن‌اند.

رابطهٔ $\Gamma \Proves !A$ بر پایهٔ استنتاج طبیعی برقرار است اگر و تنها اگر
!!^a{derivation}ی وجود داشته باشد که در آن $!A$~آخرین !!{sentence} درخت باشد
و هر برگِ !!{undischarged} در~$\Gamma$ قرار گیرد. $!A$~در استنتاج طبیعی قضیه
است اگر و تنها اگر !!^a{derivation}ی وجود داشته باشد که در آن $!A$~آخرین
!!{sentence} باشد و همهٔ فرض‌ها !!{discharged} شده باشند. برای نمونه،
!!^a{derivation} زیر نشان می‌دهد که $\Proves (!A
\land !B) \lif !A$:
\begin{prooftree}
  \AxiomC{$\Discharge{!A \land !B}{1}$}
  \RightLabel{\Elim{\land}}
  \UnaryInfC{$!A$}
  \DischargeRule{\Intro{\lif}}{1}
  \UnaryInfC{$(!A \land !B) \lif !A$}
\end{prooftree}
برچسب~$1$ نشان می‌دهد که فرض $!A \land !B$ در استنتاج \Intro{\lif}
!!{discharged} است.

مجموعهٔ~$\Gamma$ ناسازگار است اگر و تنها اگر در استنتاج طبیعی
$\Gamma \Proves \lfalse$. قاعدهٔ \FalseInt{} سبب می‌شود که از یک مجموعهٔ
ناسازگار بتوان هر !!{sentence} را !!{derive} کرد.

دستگاه‌های استنتاج طبیعی را گرهارد گنتسن و استانیسواف یاشکوفسکی در دههٔ
۱۹۳۰ پدید آوردند و سپس داگ پراویتس و فردریک فیچ آن‌ها را گسترش دادند.
از آنجا که استنتاج‌های این دستگاه‌ها روش‌های طبیعی برهان را بازتاب می‌دهند،
فیلسوفان از آن‌ها استقبال کرده‌اند. نسخه‌هایی که فیچ پروراند اغلب در
کتاب‌های درسی مقدماتی منطق به کار می‌روند. در فلسفهٔ منطق، گاه قواعد
استنتاج طبیعی را معرف معنای عملگرهای منطقی دانسته‌اند
(«معناشناسی برهان‌نظری»).

\end{document}
